
\francais
\chapter*{Résumé}

La spectroscopie Raman sert couramment à caractériser le désordre dans le graphène, mais, à notre connaissance, aucun spectre Raman résonant d'un matériau avec des défauts n'a été calculé entièrement à partir des premiers principes : les modèles existants paramètrent le couplage électron-défaut. Ce mémoire calcule, pour une lacune dans le graphène, les ingrédients \textit{ab initio} des amplitudes Raman doublement résonante, sauf le couplage électron-photon. Les intensités elles-même sont laissées à des travaux ultérieurs.

Après avoir formulé la théorie de la diffusion Raman résonante par diagrammes de Feynman, nous développons la contribution centrale de ce mémoire : la calcul du couplage électron-défaut dans une base d'ondes planes à partir de super-cellules traitées par la théorie de la fonctionnelle de la densité. Un remplissage par zéros du potentiel de défaut découple la taille de la super-cellule de la grille des fonctions de Wannier et permet l'interpolation du couplage par ces dernières sur des grilles arbitrairement denses. Le formalisme de la matrice $T$ permet de traiter la diffusion induite par le défaut à tous les ordres dans son potentiel, et la localisation du couplage dans la base de Wannier permet de résoudre cette matrice dans un amas de faible dimension : l'inversion passant de la dimension $288\,000$ à $145$ sur une grille $240\times 240$. Pour une lacune non relaxée, le couplage est dominé par sa partie non locale, le taux d'amortissement électronique dû aux défauts révèle un état quasi-lié sous la singularité de van Hove de la bande $\pi$, et l'approximation de Born surestime ce taux d'un facteur $3.3$ en médiane. Enfin, le taux d'amortissement dû aux phonons, calculé avec \texttt{EPW} dans le même sous-espace de Wannier, est du même ordre de grandeur que celui dû aux défauts à $300$~K et une concentration de $1~\%$ de lacunes.

{\bf Mots-cl\'es :} graphène, spectroscopie Raman résonante, théorie de la fonctionnelle de la densité, couplage électron-phonon, couplage électron-défaut, fonctions de Wannier, théorie de la fonctionnelle de la densité perturbative, calculs \textit{ab initio}, physique de la matière condensée.